\documentclass[]{article}
\usepackage{lmodern}
\usepackage{amssymb,amsmath}
\usepackage{ifxetex,ifluatex}
\usepackage{fixltx2e} % provides \textsubscript
\ifnum 0\ifxetex 1\fi\ifluatex 1\fi=0 % if pdftex
  \usepackage[T1]{fontenc}
  \usepackage[utf8]{inputenc}
\else % if luatex or xelatex
  \ifxetex
    \usepackage{mathspec}
  \else
    \usepackage{fontspec}
  \fi
  \defaultfontfeatures{Ligatures=TeX,Scale=MatchLowercase}
\fi
% use upquote if available, for straight quotes in verbatim environments
\IfFileExists{upquote.sty}{\usepackage{upquote}}{}
% use microtype if available
\IfFileExists{microtype.sty}{%
\usepackage{microtype}
\UseMicrotypeSet[protrusion]{basicmath} % disable protrusion for tt fonts
}{}
\usepackage[margin=1in]{geometry}
\usepackage{hyperref}
\hypersetup{unicode=true,
            pdftitle={Inching past Probabilies or Podcast Survey Results},
            pdfauthor={Martin Gleason, MS},
            pdfborder={0 0 0},
            breaklinks=true}
\urlstyle{same}  % don't use monospace font for urls
\usepackage{color}
\usepackage{fancyvrb}
\newcommand{\VerbBar}{|}
\newcommand{\VERB}{\Verb[commandchars=\\\{\}]}
\DefineVerbatimEnvironment{Highlighting}{Verbatim}{commandchars=\\\{\}}
% Add ',fontsize=\small' for more characters per line
\usepackage{framed}
\definecolor{shadecolor}{RGB}{248,248,248}
\newenvironment{Shaded}{\begin{snugshade}}{\end{snugshade}}
\newcommand{\KeywordTok}[1]{\textcolor[rgb]{0.13,0.29,0.53}{\textbf{#1}}}
\newcommand{\DataTypeTok}[1]{\textcolor[rgb]{0.13,0.29,0.53}{#1}}
\newcommand{\DecValTok}[1]{\textcolor[rgb]{0.00,0.00,0.81}{#1}}
\newcommand{\BaseNTok}[1]{\textcolor[rgb]{0.00,0.00,0.81}{#1}}
\newcommand{\FloatTok}[1]{\textcolor[rgb]{0.00,0.00,0.81}{#1}}
\newcommand{\ConstantTok}[1]{\textcolor[rgb]{0.00,0.00,0.00}{#1}}
\newcommand{\CharTok}[1]{\textcolor[rgb]{0.31,0.60,0.02}{#1}}
\newcommand{\SpecialCharTok}[1]{\textcolor[rgb]{0.00,0.00,0.00}{#1}}
\newcommand{\StringTok}[1]{\textcolor[rgb]{0.31,0.60,0.02}{#1}}
\newcommand{\VerbatimStringTok}[1]{\textcolor[rgb]{0.31,0.60,0.02}{#1}}
\newcommand{\SpecialStringTok}[1]{\textcolor[rgb]{0.31,0.60,0.02}{#1}}
\newcommand{\ImportTok}[1]{#1}
\newcommand{\CommentTok}[1]{\textcolor[rgb]{0.56,0.35,0.01}{\textit{#1}}}
\newcommand{\DocumentationTok}[1]{\textcolor[rgb]{0.56,0.35,0.01}{\textbf{\textit{#1}}}}
\newcommand{\AnnotationTok}[1]{\textcolor[rgb]{0.56,0.35,0.01}{\textbf{\textit{#1}}}}
\newcommand{\CommentVarTok}[1]{\textcolor[rgb]{0.56,0.35,0.01}{\textbf{\textit{#1}}}}
\newcommand{\OtherTok}[1]{\textcolor[rgb]{0.56,0.35,0.01}{#1}}
\newcommand{\FunctionTok}[1]{\textcolor[rgb]{0.00,0.00,0.00}{#1}}
\newcommand{\VariableTok}[1]{\textcolor[rgb]{0.00,0.00,0.00}{#1}}
\newcommand{\ControlFlowTok}[1]{\textcolor[rgb]{0.13,0.29,0.53}{\textbf{#1}}}
\newcommand{\OperatorTok}[1]{\textcolor[rgb]{0.81,0.36,0.00}{\textbf{#1}}}
\newcommand{\BuiltInTok}[1]{#1}
\newcommand{\ExtensionTok}[1]{#1}
\newcommand{\PreprocessorTok}[1]{\textcolor[rgb]{0.56,0.35,0.01}{\textit{#1}}}
\newcommand{\AttributeTok}[1]{\textcolor[rgb]{0.77,0.63,0.00}{#1}}
\newcommand{\RegionMarkerTok}[1]{#1}
\newcommand{\InformationTok}[1]{\textcolor[rgb]{0.56,0.35,0.01}{\textbf{\textit{#1}}}}
\newcommand{\WarningTok}[1]{\textcolor[rgb]{0.56,0.35,0.01}{\textbf{\textit{#1}}}}
\newcommand{\AlertTok}[1]{\textcolor[rgb]{0.94,0.16,0.16}{#1}}
\newcommand{\ErrorTok}[1]{\textcolor[rgb]{0.64,0.00,0.00}{\textbf{#1}}}
\newcommand{\NormalTok}[1]{#1}
\usepackage{graphicx,grffile}
\makeatletter
\def\maxwidth{\ifdim\Gin@nat@width>\linewidth\linewidth\else\Gin@nat@width\fi}
\def\maxheight{\ifdim\Gin@nat@height>\textheight\textheight\else\Gin@nat@height\fi}
\makeatother
% Scale images if necessary, so that they will not overflow the page
% margins by default, and it is still possible to overwrite the defaults
% using explicit options in \includegraphics[width, height, ...]{}
\setkeys{Gin}{width=\maxwidth,height=\maxheight,keepaspectratio}
\IfFileExists{parskip.sty}{%
\usepackage{parskip}
}{% else
\setlength{\parindent}{0pt}
\setlength{\parskip}{6pt plus 2pt minus 1pt}
}
\setlength{\emergencystretch}{3em}  % prevent overfull lines
\providecommand{\tightlist}{%
  \setlength{\itemsep}{0pt}\setlength{\parskip}{0pt}}
\setcounter{secnumdepth}{0}
% Redefines (sub)paragraphs to behave more like sections
\ifx\paragraph\undefined\else
\let\oldparagraph\paragraph
\renewcommand{\paragraph}[1]{\oldparagraph{#1}\mbox{}}
\fi
\ifx\subparagraph\undefined\else
\let\oldsubparagraph\subparagraph
\renewcommand{\subparagraph}[1]{\oldsubparagraph{#1}\mbox{}}
\fi

%%% Use protect on footnotes to avoid problems with footnotes in titles
\let\rmarkdownfootnote\footnote%
\def\footnote{\protect\rmarkdownfootnote}

%%% Change title format to be more compact
\usepackage{titling}

% Create subtitle command for use in maketitle
\newcommand{\subtitle}[1]{
  \posttitle{
    \begin{center}\large#1\end{center}
    }
}

\setlength{\droptitle}{-2em}
  \title{Inching past Probabilies or Podcast Survey Results}
  \pretitle{\vspace{\droptitle}\centering\huge}
  \posttitle{\par}
  \author{Martin Gleason, MS}
  \preauthor{\centering\large\emph}
  \postauthor{\par}
  \predate{\centering\large\emph}
  \postdate{\par}
  \date{2018-04-17}

\usepackage{booktabs}
\usepackage{longtable}
\usepackage{array}
\usepackage{multirow}
\usepackage[table]{xcolor}
\usepackage{wrapfig}
\usepackage{float}
\usepackage{colortbl}
\usepackage{pdflscape}
\usepackage{tabu}
\usepackage{threeparttable}
\usepackage{threeparttablex}
\usepackage[normalem]{ulem}
\usepackage{makecell}

\begin{document}
\maketitle

Probablity is not (yet) my jam. Left to my own devices, I would simply
take tutorial after tutorial, class after class, and read book after
book. Thankfully, I have an amazing friend who told me to start going to
\href{https://www.meetup.com/data-science-dojo-chicago/events/247965706/}{Meetups}.
So I attended a tutorial hosted by
\href{https://datasciencedojo.com}{Raja Iqbal}, and learned one key
lesson: Show your work.

The work, in this case, is demonstrating how I pulled survey results
from Google, cleaned it, and started visualizing it. While this project
is not directly related to developing machine learning tools for risk
classification, it has given me a better undresetanding of the process
of loading, cleaning, and visualizing real world (and useful) data.

\section{The Survey}\label{the-survey}

In preperation for launching a crowd-fund solution, my
\href{https://allcomicsconsidered.com}{podcast} crew and I decided to
ask our fans what we should do to improve our show. So we put a
\href{https://goo.gl/forms/LXkXXiJMvdSi0fpl2}{form} together and asked
the following questions:

\begin{verbatim}
##  [1] "Email Address"                                                                             
##  [2] "Why do you listen to the show? Pick all the reasons that apply."                           
##  [3] "What do you like us to talk about? Pick all the topics you like"                           
##  [4] "How did you hear about All Comics Considered?"                                             
##  [5] "How do you listen to All Comics Considered?"                                               
##  [6] "Is there another podcast platform that you'd prefer?"                                      
##  [7] "What other podcasts do you listen to?"                                                     
##  [8] "Do you listen to podcasts on YouTube?"                                                     
##  [9] "What other podcast or comic book based communities do you participate in?"                 
## [10] "What makes you want to turn All Comics Considered off?"                                    
## [11] "How would you describe All Comics Considered to someone else?"                             
## [12] "Do you suggest All Comics Considered to friends?  Why or why not?"                         
## [13] "Have you read the All Comics Considered blog in the past?"                                 
## [14] "Would you be interested in more written content?"                                          
## [15] "How much do you like the Pullbox Feature?"                                                 
## [16] "What would improve the Pullbox Feature?"                                                   
## [17] "How much do you enjoy our deep dive features into specific comics, properties, and issues?"
## [18] "How much do you enjoy our shorter, more comedic segments?"                                 
## [19] "How do you feel about the length of the show?"                                             
## [20] "What format changes might you be interested in seeing come to the show?"
\end{verbatim}

26 people responded to these question. Google Forms has some simple
analytics built into the platform, and given the number of responses, a
good deal of the insights we gathered from our survey was pulled from
that front page. But pulling insights from a form was not how I wanted
to access the responses. So I took the opportunity this survey created
to see how well I could use R to pull the data.

\section{Making it More Complicated}\label{making-it-more-complicated}

The first step was determining which package I should use to pull the
responses. It took a grand total of two minutes to decide on the
package: \texttt{googlesheets}. This
\href{https://cran.r-project.org/web/packages/googlesheets/index.html}{package},
by \href{https://github.com/jennybc/googlesheets}{Jennifer Bryan} and
Joanna Zhao, has simple, clear, and consistent syntax. Once my
\href{https://rawgit.com/jennybc/googlesheets/master/vignettes/managing-auth-tokens.html}{Google
credentials}\footnote{Jennifer's vingette explains authentication, where
  tokens live, and everything you wanted to know about how the package
  works.} were authenticated, pulling data from the sheets was a breeze.
The function \texttt{gs\_ls()} pulled a list of my sheets into a tibble,
which made finding the key to the survey responses as simple as:
\texttt{gs\_ls\ \%\textgreater{}\%\ filter(str\_detect(sheet\_title,\ "All\ Comics\ Considered:\ 2018"))\ \%\textgreater{}\%\ select(sheet\_key)\ \%\textgreater{}\%\ as.character()}.\footnote{I
  was not, originally, so clever with my code. First I stored my sheets
  in a tibble like so, \texttt{sheets\ \textless{}-\ gs\_ls()}, then I
  may have cut and pasted instead of using \texttt{stringr} and
  \texttt{as.character()}. Writing this article now, I realize it could
  have been as easy as
  \texttt{gs\_ls()\ \%\textgreater{}\%\ filter(str\_detect(sheet\_title,\ "All\ Comics\ Considered:\ 2018"))\ \%\textgreater{}\%\ select(sheet\_key)\ \%\textgreater{}\%\ as.character()\ \%\textgreater{}\%\ gs\_key()\ \%\textgreater{}\%\ gs\_read()}.}

Armed with the key, I pulled the survey resultls into a single tibble:

\begin{Shaded}
\begin{Highlighting}[]
\NormalTok{register_survey <-}\StringTok{ }\KeywordTok{gs_key}\NormalTok{(survey_key)}
\NormalTok{survey_full <-}\StringTok{ }\KeywordTok{gs_read}\NormalTok{(register_survey)}
\NormalTok{survey_anon <-}\StringTok{ }\NormalTok{survey_full }\OperatorTok{%>%}\StringTok{ }\KeywordTok{select}\NormalTok{(}\OperatorTok{-}\DecValTok{2}\NormalTok{)}
\end{Highlighting}
\end{Shaded}

This could have been collapsed into one line of code, which I write in
footnote 2. Both methods produce the following table

Overwhelming Table

Timestamp

Why do you listen to the show? Pick all the reasons that apply.

What do you like us to talk about? Pick all the topics you like

How did you hear about All Comics Considered?

How do you listen to All Comics Considered?

Is there another podcast platform that you'd prefer?

What other podcasts do you listen to?

Do you listen to podcasts on YouTube?

What other podcast or comic book based communities do you participate
in?

What makes you want to turn All Comics Considered off?

How would you describe All Comics Considered to someone else?

Do you suggest All Comics Considered to friends? Why or why not?

Have you read the All Comics Considered blog in the past?

Would you be interested in more written content?

How much do you like the Pullbox Feature?

What would improve the Pullbox Feature?

How much do you enjoy our deep dive features into specific comics,
properties, and issues?

How much do you enjoy our shorter, more comedic segments?

How do you feel about the length of the show?

What format changes might you be interested in seeing come to the show?

3/3/2018 20:14:32

Timely and relevant content., Social relevance., The selection of books
talked about.

Comics TV, Comic Movies, Historical comic titles, Social Justice Issues
in Comics

facebook

Google Play

itunes

NA

Yes

NA

NA

itchy

sure

Yes

Yes

5

NA

5

5

3

Full episodes about a single topic

3/3/2018 20:30:15

Timely and relevant content., Banter between hosts., Social relevance.,
The selection of books talked about., Interviews., The best show I have
found for balancing love of the medium with a commitment to keeping its
feet to the fire when it comes to representation.

Comics TV, Comic Movies, Ongoing comic titles, Historical comic titles,
Social Justice Issues in Comics, Ourselves, I think you have the balance
right now. Being album to compare the scenarios in different comics to
aspects of your own life makes your analysis more grounded.

Websearch. Serendipity.

iTunes

I use Otto occasionally. Mainly iTunes

Where to start: 95BFM; Adventure: A dungeons and dragons podcast; All in
the mind; All songs considered; RNZ At the movies; Book Riot; Can He Do
That?; Talking Comic; Freakonomics Radio; Fresh Air; Gone by Lunchtime;
Larry Wilmore; Mental Illness Happy Hour; NPR Politics Podcast; On One
with Angela Rye; Pod Save America; Pod Save the World; Pop Culture Happy
Hour; Race Invaders; Still ProcessingThe Ezra Klein Show; Up, Up and
Away; WTF with Marc Maron; Women of Marvel

No

I listen to a lot of Podcasts, but this is the only one I really
interact with on Facebook etc.

Nothing.

Smart, witty, great coverage of the intersection of politics and comics.

I've actually included the podcast in a `recommended' list on a radio
spot I used to do for my local student radio station.

NA

Yes

5

I like the idea of regular Deep Cuts (great story arcs from the past)

5

4

2

Full episodes about a single topic

This table violates one of the most critical lessons I learned in grad
school: The values are not atomic. In other words, the answer to the
question \texttt{names(survey\_anon{[}2{]})} may have more than one
answer per respondant:

\begin{verbatim}
##  [1] "Timely and relevant content., Social relevance., The selection of books talked about."                                                                                                                                                                                    
##  [2] "Timely and relevant content., Banter between hosts., Social relevance., The selection of books talked about., Interviews., The best show I have found for balancing love of the medium with a commitment to keeping its feet to the fire when it comes to representation."
##  [3] "Banter between hosts., Social relevance., The selection of books talked about., Interviews."                                                                                                                                                                              
##  [4] "Banter between hosts., Social relevance., Interviews."                                                                                                                                                                                                                    
##  [5] "Banter between hosts., Social relevance., The selection of books talked about."                                                                                                                                                                                           
##  [6] "Banter between hosts., The selection of books talked about."                                                                                                                                                                                                              
##  [7] "Timely and relevant content., Banter between hosts., Social relevance."                                                                                                                                                                                                   
##  [8] "Banter between hosts., Social relevance."                                                                                                                                                                                                                                 
##  [9] "Timely and relevant content., Banter between hosts., Social relevance., Interviews."                                                                                                                                                                                      
## [10] "Timely and relevant content., Banter between hosts."                                                                                                                                                                                                                      
## [11] "Timely and relevant content., Banter between hosts., Social relevance., Interviews., Perspective and theory"                                                                                                                                                              
## [12] "Banter between hosts., Social relevance."                                                                                                                                                                                                                                 
## [13] "Timely and relevant content., Banter between hosts., Social relevance., The selection of books talked about., Interviews., Pretending that your office is a fun place full of friends is nice."                                                                           
## [14] "Timely and relevant content., Social relevance."                                                                                                                                                                                                                          
## [15] "Banter between hosts., Social relevance., Interviews."                                                                                                                                                                                                                    
## [16] "Banter between hosts., The selection of books talked about., Interviews."                                                                                                                                                                                                 
## [17] "The selection of books talked about., Interviews."                                                                                                                                                                                                                        
## [18] "Timely and relevant content., Banter between hosts."                                                                                                                                                                                                                      
## [19] "Timely and relevant content., Banter between hosts., Social relevance., The selection of books talked about., Interviews."                                                                                                                                                
## [20] "Banter between hosts., Social relevance., The selection of books talked about., To privately laugh at Marty without hurting his feelings."                                                                                                                                
## [21] "Timely and relevant content., Banter between hosts., Interviews."                                                                                                                                                                                                         
## [22] "Timely and relevant content., Social relevance."                                                                                                                                                                                                                          
## [23] "Timely and relevant content., The selection of books talked about."                                                                                                                                                                                                       
## [24] "Timely and relevant content., Social relevance."                                                                                                                                                                                                                          
## [25] "The selection of books talked about."                                                                                                                                                                                                                                     
## [26] "Timely and relevant content., Social relevance., The selection of books talked about."
\end{verbatim}


\end{document}
